\section{Setting}
\label{sec:setting}

\paragraph{Models.} Holo-3.1 models are fine-tunes of the Qwen3.5 vision-language family \citep{hcompany2026holo31}. Their decoder is hybrid, with three linear-attention (gated DeltaNet) layers for every full-attention layer and tied input/output embeddings. Holo-3.1-4B has \cfgBigLayers{} layers of width \cfgBigHidden{} (MLP width \cfgBigInter); Holo-3.1-0.8B has \cfgSmallLayers{} layers of width \cfgSmallHidden. Layers are indexed from zero.

\paragraph{OSWorld harness (4B).} We reuse the Holo agent loop as packaged by RealtimeGUI at commit \texttt{a763361} \citep{realtimegui}, a third-party reproduction rather than an H Company repository, changing only the model provider. Its rule for recognising a declared infeasibility (\texttt{task.\{0,4\}infeasible|infeasible|cannot be (completed|done)|not possible}) is broader than that of the official OSWorld adapter for Holo, which accepts only \texttt{task.\{0,4\}infeasible}; our absolute scores are therefore not strictly comparable with H Company's, while paired comparisons within our harness are. % source: data/local/harness/run_holo.py (_INFEASIBLE_RE), data/local/harness/surfer_agent.py (official xlang-ai/OSWorld mm_agents/surferH); git remote of ~/osw/rtg = github.com/1azybug/RealtimeGUI @a763361; data/local/JOURNAL.md 02:25 UTC
vLLM~0.23 serves the 4B with thinking enabled, $T=0.8$, at most 50 steps and the three latest screenshots in context. % source: JOURNAL.md "OSWorld — infra prête" and "OSWorld — pilote"
OSWorld \citep{xie2024osworld} tasks run in Ubuntu VMs under Docker+KVM, at most four in parallel, and are scored by the benchmark's execution-based checkers. We fixed three infrastructure issues in the VM provider (a memory cap, binding VM ports to localhost, an nftables network script) and modified no task or checker. % source: JOURNAL.md "3 correctifs d'infra du provider docker"

\paragraph{Frozen split.} From the 361 tasks of RealtimeGUI's \texttt{test\_nogdrive.json}, we sort ids by SHA-1 within each domain and deal them alternately into a \textsc{diag} and a \textsc{test} pool, keeping 120 tasks each with per-domain quotas. The split files have SHA-1 prefixes \texttt{51e670b3} (diag) and \texttt{957d68e7} (test). No construction or screening read the test split. Four online runs then used it (\S\ref{sec:online}): the original, the ExPO edit, which went to it without having passed the gate, a partial corrective-prompt run, and the CUA-Lite merge under a decision rule fixed in advance. % source: JOURNAL.md "OSWorld — split figé"; data/local/JOURNAL.md 03:05 UTC; data/local/reports/REDTEAM_0305.md point 5

\paragraph{MiniWoB++ harness (0.8B).} Our first attempt used Holo-3.1-0.8B on MiniWoB++ \citep{shi2017wob,liu2018miniwob} through BrowserGym \citep{browsergym2024}, with the task area cropped and upscaled $5\times$; it is described in \S\ref{sec:mw} and Appendix~\ref{app:mw}.

\paragraph{Hardware and licences.} Rollouts, replays and screens ran on one NVIDIA H200 shared between the vLLM server and research jobs; edited checkpoints were also served on rented H100 GPUs, and the OSWorld VMs ran on a separate CPU host. As order-of-magnitude estimates, the project used about 14 hours of one H200 (rented, shared by every research job and the rollout server, so per-experiment GPU-hours are not separable), about 8 H100 GPU-hours (USD 31.9 billed over two accounts) and about 60 GitHub Actions CPU jobs for the OSWorld VMs. % source: data/local/MODAL_COST.md (16.65 + 15.26 USD at 05:11 UTC, H100 at about 4 USD/h); H200 hours and job count: coordinator estimates (ares, 05:25 UTC)
% source: data/local/JOURNAL.md ("vLLM 0.23 sur H200", 01:10 UTC "Modal H100"); data/pod_astra/merge/VERIFY_CUA_B050_BEHAVIOR.md (H200 vs H100)
Holo-3.1-0.8B and -4B, Qwen3.5-4B, the CUA-Lite donor, OSWorld and BrowserGym are released under Apache-2.0 and MiniWoB++ under MIT; the RealtimeGUI repository has no licence file at the cited commit, and we use it unmodified except for the model provider. % source: Hugging Face model API and GitHub license API, queried 2026-11-08
